\documentclass[english,11pt,a4paper]{article}
\usepackage[margin=1in]{geometry}
\usepackage{fontspec}
\setmainfont{lmroman10-regular.otf}[BoldFont=lmroman10-bold.otf,ItalicFont=lmroman10-italic.otf,BoldItalicFont=lmroman10-bolditalic.otf]
\setsansfont{lmsans10-regular.otf}[BoldFont=lmsans10-bold.otf]
\setmonofont{lmmono10-regular.otf}[Scale=MatchLowercase]
\newfontfamily\LeanSymbols{seguisym.ttf}[Path=typeset_runtime/]
\newlength{\LeanCharWidth}
\usepackage{amsmath,amssymb}
\usepackage[english]{babel}
\usepackage{microtype,parskip}
\usepackage{graphicx,float,caption}
\usepackage{longtable,booktabs,array,calc}
\usepackage{xcolor}
\usepackage{xurl}
\usepackage{hyperref,bookmark}
\usepackage{etoolbox,needspace}
\setcounter{secnumdepth}{-\maxdimen}
\setlength{\emergencystretch}{3em}
\providecommand{\tightlist}{\setlength{\itemsep}{0pt}\setlength{\parskip}{0pt}}
\newcounter{none}
\AtBeginEnvironment{longtable}{\Needspace{8\baselineskip}\small}
\pretocmd{\section}{\Needspace{5\baselineskip}}{}{}
\pretocmd{\subsection}{\Needspace{4\baselineskip}}{}{}
\pretocmd{\subsubsection}{\Needspace{4\baselineskip}}{}{}
\captionsetup{font=small,labelfont=bf}
\urlstyle{same}
\providecommand{\pandocbounded}[1]{#1}
\hypersetup{pdftitle={Paper II: formalization and reproducibility},pdfauthor={Juan Pablo Traverso Gianini},pdflang={en-US},colorlinks=true,linkcolor=blue,urlcolor=blue,citecolor=blue}
\title{Paper II: formalization and reproducibility}
\author{Juan Pablo Traverso Gianini\\\small Independent researcher, Santiago, Chile\\\small \href{mailto:jtraverso@gmail.com}{jtraverso@gmail.com}\\\small ORCID: 0009-0003-6068-4096}
\date{}
\begin{document}
\raggedbottom
\maketitle
\textbf{Paper II · Formal companion v1.3}\\
19 October 2026

This document accompanies the mathematical article. It preserves the
formal declarations, scope qualifications and computational evidence
previously included in the article. The source freeze and its recorded
verification are unchanged. No Lean/Lake build was run for this edition.
Historical section and table numbers in the records below refer to the
preceding article; they are retained as locators, not renumbered as new
verification results. Original references {[}1{]}, {[}2{]}, etc. are
listed at the end.

\href{https://github.com/jtraverso/erdos-81-chordal-clique-partitions/tree/eb335845bfa23375d5040ea31fc4f4b366346287/preprints/PAPER_II/05_formalization}{Frozen
sources and historical evidence / Fuentes congeladas y evidencias
históricas} ·
\href{https://github.com/jtraverso/erdos-81-chordal-clique-partitions/tree/eb335845bfa23375d5040ea31fc4f4b366346287/preprints/PAPER_II/01_manuscript}{Previous
article / Artículo anterior}

\subsubsection{\texorpdfstring{Corollary 1.2\(^{\prime}\) (sharpened
asymptotic)}{Corollary 1.2\^{}\{\textbackslash prime\} (sharpened asymptotic)}}\label{corollary-1.2prime-sharpened-asymptotic}

\textbf{Formalization status.} The delivered Paper II Lean v1.2 freeze
uses Lean 4 / Mathlib v4.28.0 {[}5,6{]}. Its recorded main build
completed 8,061 jobs, and the explicit extremizer/copy-defect supplement
completed 8,032 jobs. The main declaration is named
\nolinkurl{PaperII.theorem_1_2}; it corresponds to Theorem 1.1 of this
manuscript and includes the accompanying attainment statement. The same
freeze includes the complete-split argmax, level sets, copy defects, and
\nolinkurl{PaperII.AsymptoticCorollaries}. Its theorem-level axiom gate
reports no project-level mathematical axiom. The formal development
derives the chordal structure used in the proof from the standard
predicate \nolinkurl{IsChordal}; no clique-tree or chordal-structure axiom
is added. Section 8 records the exact freeze perimeter.

\subsection{8. Reproducibility, formal verification, and
scope}\label{reproducibility-formal-verification-and-scope}

The proof in the paper is analytic. Computational experiments were used
during development to search for counterexamples and to catch
transcription errors, but no computation is used as a premise in the
proof.

The accompanying calculation ledger records the proof at a lower level
of granularity. It lists the definitions, the copied-cover construction
in Lemma 3.1, the two cancellation identities, the discrete-convexity
lift to clone classes, the clone-progress argument, the simpliciality
requirement in the chordal reduction, the terminal complete-split
characterization, the two-variable program on \(S_{p,q}\), the
degenerate case \(S_{2,0}\), the formulas for \(\Phi_\tau(S_{p,q})\),
and the final integer maximization. The theorem, the exact bound, and
the assembly have received internal adversarial audit. This records the
author's validation process and is not a claim of external peer review.

\subsubsection{Formal verification in
Lean}\label{formal-verification-in-lean}

The consolidated snapshot at
\nolinkurl{PAPER_II/05_formalization/lean_v1.2_freeze} uses Lean 4 /
Mathlib v4.28.0 {[}5,6{]}. The recorded main build completed 8,061 jobs
with zero errors; an explicit supplement covering
\nolinkurl{PaperII.Extremizer} and \nolinkurl{PaperII.CopyDefect} completed
8,032 jobs with zero errors. The corresponding main declaration is

\[
\texttt{PaperII.theorem\_1\_2}
\]

in repository {[}7{]}. The formal theorem includes both sides of the
result: the upper bound over all chordal graphs on \(n\) vertices and
the existence of a complete-split graph attaining the value.

The formalization uses the standard chordal predicate
\nolinkurl{IsChordal}. In the development, this predicate says that every
cycle of length at least four has a chord. The deficit functional is
formalized as

\[
\text{\LeanSymbols phiTau G = (|E(G)| : ℝ) - 2 · tau3star G},
\]

where \nolinkurl{tau3star} is the fractional triangle-cover number of the
finite cover LP. The complete-split family is formalized as
\nolinkurl{completeSplit\ p\ q}, and the extremal value is the integer
\(\lfloor(2n+1)^2/24\rfloor\).

No chordal-structure axiom is added. The formalization derives from
\nolinkurl{IsChordal} the hereditary property for induced subgraphs,
preservation under adding a vertex with clique neighborhood, and the
existence of simplicial vertices. The written proof of Lemma 5.1 uses
clique trees as standard textbook infrastructure; the Lean proof uses a
clique-tree-free terminal-characterization argument. Thus the formal
proof does not assume a clique-tree or running-intersection axiom.

The checked development covers the copied-cover inequality, clone-class
discrete convexity, clone-class progress, chordality-preserving copying,
the terminal complete-split characterization, orbit averaging on
\(S_{p,q}\), the terminal linear program, and the exact integer
maximization. The freeze log records 49 linter warnings, principally
unused section variables and unused simp arguments; these do not affect
kernel elaboration. The recorded axiom report is

\begin{quote}
{\LeanSymbols\fontsize{8}{10}\selectfont
\noindent \nolinkurl{#print}\hspace{0.45em}\nolinkurl{axioms}\hspace{0.45em}\nolinkurl{PaperII.theorem_1_2}\par
\noindent \nolinkurl{[propext,}\hspace{0.45em}\nolinkurl{Classical.choice,}\hspace{0.45em}\nolinkurl{Quot.sound]}\par
}
\end{quote}

These are Lean's standard classical axioms. The report contains no
\nolinkurl{sorryAx} and no project-specific axiom. The frozen source
archive is \nolinkurl{PAPER_II_lean_v1.2_freeze.zip}, with SHA-256
\nolinkurl{ee2d05cc40d943ca92f8f7bf3e5dd83c2692518ddea5e2ca4f7686ccb1ac3895}.
The main and supplement logs are stored as
\nolinkurl{gate_logs/BUILD_LOG.txt} and
\nolinkurl{gate_logs/BUILD_SUPPLEMENT_EXTREMIZER_COPYDEFECT.txt} within
\nolinkurl{lean_v1.2_freeze}. The chordal-graph contribution modules
shipped in the same archive are candidates for upstreaming to Mathlib
and are not additional proof assumptions of this paper. The frozen
archive passed independent clean-room reproduction; the external audit
recorded a successful 8,063-job build and the expected 16 axiom
surfaces.

\textbf{Table 4. Frozen formalization perimeter (Lean v1.2).} The
delivered freeze contains the main theorem and the manuscript
declarations listed below; independent clean-room reproduction returned
PASS.

{\def\LTcaptype{none} % do not increment counter
\begin{longtable}[]{@{}
  >{\raggedright\arraybackslash}p{(\linewidth - 4\tabcolsep) * \real{0.3333}}
  >{\raggedright\arraybackslash}p{(\linewidth - 4\tabcolsep) * \real{0.3333}}
  >{\raggedright\arraybackslash}p{(\linewidth - 4\tabcolsep) * \real{0.3333}}@{}}
\toprule\noalign{}
\begin{minipage}[b]{\linewidth}\raggedright
Node
\end{minipage} & \begin{minipage}[b]{\linewidth}\raggedright
Statement
\end{minipage} & \begin{minipage}[b]{\linewidth}\raggedright
Axiom footprint
\end{minipage} \\
\midrule\noalign{}
\endhead
\bottomrule\noalign{}
\endlastfoot
\nolinkurl{PaperII.theorem_1_2} &
\(\max_{\text{chordal},\,|V|=n}\Phi_\tau=\lfloor(2n+1)^2/24\rfloor\),
attained by a complete-split graph & \nolinkurl{propext},
\nolinkurl{Classical.choice}, \nolinkurl{Quot.sound} \\
\nolinkurl{Fsat_argmax_unique}, \nolinkurl{Fsat_argmax_tie} & uniqueness
of the maximizing integer \(p\) within the complete-split family, and
the two-value tie when \(n\equiv1\pmod3\) & \nolinkurl{propext},
\nolinkurl{Classical.choice}, \nolinkurl{Quot.sound} \\
\nolinkurl{level_set_iff} (Cor. 7.2) & exact level-set equivalence from
the continuous maximum & \nolinkurl{propext}, \nolinkurl{Classical.choice},
\nolinkurl{Quot.sound} \\
\nolinkurl{copyDefect_nonneg}, \nolinkurl{copyGamma_ge_half_copyDefect} &
copy defect \(\Delta_{\Phi_\tau}\ge0\) and
\(\Gamma_{\Phi_\tau}\ge\Delta_{\Phi_\tau}/2\) & \nolinkurl{propext},
\nolinkurl{Classical.choice}, \nolinkurl{Quot.sound} \\
\end{longtable}
}

\textbf{Table 5. Arithmetic surface included in the Lean v1.2 freeze.}
These declarations are consequences of the frozen extremal value and are
included in the delivered archive.

{\def\LTcaptype{none} % do not increment counter
\begin{longtable}[]{@{}
  >{\raggedright\arraybackslash}p{(\linewidth - 4\tabcolsep) * \real{0.3333}}
  >{\raggedright\arraybackslash}p{(\linewidth - 4\tabcolsep) * \real{0.3333}}
  >{\raggedright\arraybackslash}p{(\linewidth - 4\tabcolsep) * \real{0.3333}}@{}}
\toprule\noalign{}
\begin{minipage}[b]{\linewidth}\raggedright
Node
\end{minipage} & \begin{minipage}[b]{\linewidth}\raggedright
Statement
\end{minipage} & \begin{minipage}[b]{\linewidth}\raggedright
Axiom footprint
\end{minipage} \\
\midrule\noalign{}
\endhead
\bottomrule\noalign{}
\endlastfoot
\nolinkurl{phiTau_max_sandwich} (Cor. 1.2\(^{\prime}\)) & for every
\(n\in\mathbb Z\),
\(4n^2+4n-23<24\lfloor(2n+1)^2/24\rfloor\le4n^2+4n+1\) &
\nolinkurl{propext}, \nolinkurl{Classical.choice}, \nolinkurl{Quot.sound} \\
\nolinkurl{odd_sq_emod_24}, \nolinkurl{phiTau_max_closed} (Cor. 7.3) &
\((2n+1)^2\equiv1\) or \(9\pmod{24}\), and the resulting residue-class
formula for the extremum & \nolinkurl{propext}, \nolinkurl{Classical.choice},
\nolinkurl{Quot.sound} \\
\nolinkurl{phiTau_max_le_paperI_bound} &
\(\lfloor(2n+1)^2/24\rfloor\le n^2/6+n/2\) for \(n\ge1\) (comparison
with Paper I) & \nolinkurl{propext}, \nolinkurl{Quot.sound} \\
\end{longtable}
}

The four declarations in Table 5 are collected in
\nolinkurl{PaperII/AsymptoticCorollaries.lean}, which builds against the
same toolchain (Lean 4 / Mathlib v4.28.0) and carries only Lean's
standard classical axioms, with no \nolinkurl{sorryAx} and no
project-specific axiom. The per-declaration \nolinkurl{#print\ axioms}
reports are archived in \nolinkurl{PaperII/AxiomCheckCorollaries.lean}. The
module adds no hypothesis to and does not modify the frozen main
development; each declaration is derived from the closed value by
\nolinkurl{ring}, a finite residue split, and \nolinkurl{omega}.

\subsubsection{Independent computational
verification}\label{independent-computational-verification}

In addition to the analytic proof and the Lean development,
supplementary scripts recomputed the main claims from first principles.
These scripts evaluated \(\tau_3^*\) by exact linear programming,
checked the dual fractional packing LP separately as a consistency test,
performed the integer optimization in exact rational arithmetic, and
used exhaustive enumeration where feasible. The supplied audit report
records no discrepancy.

The reported checks include:

\begin{itemize}
\tightlist
\item
  the vertex-copy inequality and its two cancellation identities over
  all non-isomorphic graphs on up to \(7\) vertices and all nonadjacent
  pairs;
\item
  discrete convexity, clone-class progress without splitting, and
  chordality preservation, including a \(C_4\) test showing why the
  simpliciality hypothesis is necessary;
\item
  the terminal characterization over all chordal graphs on up to \(7\)
  vertices;
\item
  the complete-split formulas of Proposition 6.2 and Corollary 6.3
  against the full LP in the tested range;
\item
  the identity
  \(\max_{p+q=n}\Phi_\tau(S_{p,q})=\lfloor(2n+1)^2/24\rfloor\) for
  \(1\le n\le2000\);
\item
  the global maximum over chordal graphs, exhaustively for \(n\le7\),
  with additional random tests for \(8\le n\le12\).
\end{itemize}

The audit scripts, reports, and SHA-256 manifests accompany the
preprint. They corroborate the proof and help make the release
reproducible; they are not a substitute for the proof.

\subsubsection{Scope}\label{scope}

The result proved here is the exact finite extremum of the fractional
functional \(\Phi_\tau\). It is not a clique-partition theorem and does
not use integral triangle packing. Within the Erdős Problem \#81
program, the role of this paper is to identify the complete-split
terminal family and the exact fractional value.

\begin{center}\rule{0.5\linewidth}{0.5pt}\end{center}

\subsubsection{Corollary 7.3 (residue-class
formula)}\label{corollary-7.3-residue-class-formula}

Thus the extremal value is one of these two integers according to the
residue of \(n\) modulo \(3\). This makes explicit the
odd-square-modulo-\(24\) fact used in the case analysis of Proposition
7.1. The same residue class \(n\equiv1\pmod3\) is exactly the two-value
tie case recorded by \nolinkurl{Fsat_argmax_tie}. The residue fact is
formalized as \nolinkurl{PaperII.odd_sq_emod_24} and the closed form as
\nolinkurl{PaperII.phiTau_max_closed} (axiom footprint \nolinkurl{propext},
\nolinkurl{Classical.choice}, \nolinkurl{Quot.sound}).

\subsubsection{Lemma 5.1}\label{lemma-5.1}

\emph{The accompanying Lean formalization proves this lemma by an
equivalent argument that does not use clique trees --- via simplicial
elimination and the fact that every minimal vertex separator of a
chordal graph is a clique --- rather than the clique-tree proof given
below; see §8. The statement proved is identical.}

\subsubsection{Theorem 1.1}\label{theorem-1.1}

with \(p\) an integer nearest to \((2n+1)/6\). Among complete-split
graphs \(S_{p,n-p}\), the maximizing integer \(p\) is unique unless
\(n\equiv1\pmod3\), in which case exactly two consecutive values of
\(p\) attain the maximum; see the case analysis in the proof of
Proposition 7.1 and the formalized declarations
\nolinkurl{Fsat_argmax_unique} and \nolinkurl{Fsat_argmax_tie}. No
uniqueness statement for arbitrary chordal extremizers is claimed.

\subsubsection{\texorpdfstring{Corollary 1.2\(^{\prime}\) (sharpened
asymptotic)}{Corollary 1.2\^{}\{\textbackslash prime\} (sharpened asymptotic)}}\label{corollary-1.2prime-sharpened-asymptotic-1}

After the explicit linear term \(n/6\) is extracted, the remaining error
is bounded: one has \(M(n)=n^2/6+n/6+\theta_n\) with
\(\theta_n\in(-1,\,1/24]\). This is immediate from
\(24\,M(n)\le(2n+1)^2<24\,M(n)+24\) and \((2n+1)^2=4n^2+4n+1\).
Formalized for \(n:\mathbb Z\), with no lower-bound hypothesis, as
\nolinkurl{PaperII.phiTau_max_sandwich} (axiom footprint
\nolinkurl{propext}, \nolinkurl{Classical.choice}, \nolinkurl{Quot.sound}).

\subsubsection{Remark (coherence with Paper I's
bound)}\label{remark-coherence-with-paper-is-bound}

where the right-hand side is the finite estimate
\(|E(G)|-2\nu_3^*(G)\le n^2/6+n/2\) from Paper I {[}4{]}. This
comparison is numerical rather than a consequence of class containment,
since the chordal class is larger than the split class. The slack is
\(\bigl(n^2/6+n/2\bigr)-M(n)=n/3+O(1)\). The inequality is formalized as
\nolinkurl{PaperII.phiTau_max_le_paperI_bound} (axiom footprint
\nolinkurl{propext}, \nolinkurl{Quot.sound}).

\subsubsection{Definition (copy defect)}\label{definition-copy-defect}

For nonadjacent \(u,v\) set \[
\begin{aligned}
\Delta_{\Phi_\tau}(u,v)
&=\Phi_\tau(G_{v\to u})+\Phi_\tau(G_{u\to v})-2\Phi_\tau(G)\ge0
\quad(\text{Lemma 3.1}),\\
\Gamma_{\Phi_\tau}(u,v)
&=\max\{\Phi_\tau(G_{v\to u}),\Phi_\tau(G_{u\to v})\}-\Phi_\tau(G).
\end{aligned}
\] Then \(\Gamma_{\Phi_\tau}\ge\Delta_{\Phi_\tau}/2\), formalized as
\nolinkurl{copyDefect_nonneg} and
\nolinkurl{copyGamma_ge_half_copyDefect}. The quantity
\(\Delta_{\Phi_\tau}\) is the resource that any future stability
statement must control; it is recorded here as a quantitative handle,
not as a stability theorem.

\subsubsection{Corollary 7.2 (level sets of the cover
functional)}\label{corollary-7.2-level-sets-of-the-cover-functional}

Write \(M=2n+1\). For \(p\) in the saturated branch \(q\ge p-1\) and
every \(\delta\ge0\), \[
\Phi_\tau(S_{p,q})\ \ge\ \frac{M^2}{24}-\delta\qquad\Longleftrightarrow\qquad \Big|p-\frac{M}{6}\Big|\ \le\ \sqrt{\frac{2\delta}{3}},
\] with \(\delta\) measured from the continuous maximum \(M^2/24\). This
is an exact equivalence, obtained from identity (7.1) by monotonicity of
the square, and is formalized as \nolinkurl{level_set_iff}. In particular
\(\Phi_\tau\) is not flat in any direction of \(p\), and its level sets
are intervals of radius \(\sqrt{2\delta/3}\); within the complete-split
family this settles the second half of the stability question raised in
Paper III, §12.4.

\subsection{References}\label{references}

{[}1{]} P. Erdős, E. T. Ordman, and Y. Zalcstein, ``Clique partitions of
chordal graphs,'' \emph{Combinatorics, Probability and Computing}
\textbf{2} (1993), no. 4, 409--415.

{[}2{]} M. C. Golumbic, \emph{Algorithmic Graph Theory and Perfect
Graphs}, second edition, Annals of Discrete Mathematics 57, Elsevier,
Amsterdam, 2004.

{[}3{]} J. R. S. Blair and B. Peyton, ``An introduction to chordal
graphs and clique trees,'' in \emph{Graph Theory and Sparse Matrix
Computation}, IMA Volumes in Mathematics and its Applications 56,
Springer, New York, 1993, pp.~1--29.

{[}4{]} J. P. Traverso Gianini, ``Affine profile reduction for
fractional triangle packings in split graphs'' (Paper I in the series),
preprint, July 2026.

{[}5{]} L. de Moura and S. Ullrich, ``The Lean 4 Theorem Prover and
Programming Language,'' in \emph{Automated Deduction -- CADE 28},
Lecture Notes in Computer Science \textbf{12699}, Springer, 2021,
pp.~625--635, \url{https://doi.org/10.1007/978-3-030-79876-5_37}.

{[}6{]} The mathlib Community, ``The Lean Mathematical Library,'' in
\emph{Proceedings of the 9th ACM SIGPLAN International Conference on
Certified Programs and Proofs (CPP 2020)}, ACM, 2020, pp.~367--381,
\url{https://doi.org/10.1145/3372885.3373824}.

{[}7{]} J. P. Traverso Gianini, \emph{Erdős \#81 Chordal Clique
Partitions: Public Preprints and Formalization Artifacts}, public
project repository, 2026,
\url{https://github.com/jtraverso/erdos-81-chordal-clique-partitions},
accessed July 7, 2026.

{[}8{]} T. F. Bloom, ``Erdős Problem \#81,'' \emph{Erdős Problems},
\url{https://www.erdosproblems.com/81}, accessed July 6, 2026.

\end{document}
